$R_{0,0,0,0}$ is seen to satisfy \eqref{eqn:m10init} easily.
Due to the shifts \eqref{eqn:shiftRx} and \eqref{eqn:shiftRy}, showing that $R_{0,0,0,0}$ satisfies
\eqref{eqn:m10d1} is equivalent to showing:
\begin{align}
R_{0,0,0,0} - R_{0,2,1,1} - xqR_{0,4,2,2} = 0.
\end{align}
This relation can be obtained as a linear combination of the fundamental relations \eqref{eqn:relm1}--\eqref{eqn:relm4}:
\begin{align}
-&m_1(-2,0,0,0)+m_1(-2,0,0,1)-xq\cdot m_1(0,4,2,2)
\nonumber\\
&+xq\cdot m_1(0,4,3,2)+m_2(0,0,0,1)+m_3(0,2,0,1)
\nonumber \\
&-xq\cdot m_3(2,4,2,2)+m_4(-2,0,0,0)-y\cdot m_4(2,6,4,2).
\end{align}

Now we prove the uniqueness. 
Suppose that $F(x,y,q)=\sum_{i,j,k\geq 0}f_{i,j,k}x^iy^jq^k$ is a solution to this system.
Note that for all $i,j,k\geq 0$, $f_{0,j,k}$, $f_{i,0,k}$ and $f_{i,j,0}$ 
are uniquely determined due to the initial conditions 
\eqref{eqn:m10init}.
For convenience, we define $f_{i,j,k}=0$ whenever any one or more of $i,j,k$ are negative.
\eqref{eqn:m10d1} translates to:
\begin{align}
f_{i,j,k}&=f_{i,j,k-i}+f_{i-1,j,k-2i+1}.
\label{eqn:frel}
\end{align}
Now we induct on $N=i+k$ and show that all $f_{i,j,k}$ are uniquely determined.
The case~$N=0$ follows easily since $f_{0,j,0}$ are uniquely known due to \eqref{eqn:m10init}.
Suppose that for all triples $(i,j,k)$ with $i+k<N$, the values $f_{i,j,k}$ are determined.
Pick a triple~$(I,J,K)$ such that $I+K=N$. 
If $I=0$, we know $f_{0,J,N}$ by \eqref{eqn:m10init}.
So, suppose that~$I\geq 1$. 
Then, the RHS of \eqref{eqn:frel} involves two terms:
\begin{enumerate}
\item $f_{I,J,K-I}$ for which $I+(K-I) = K < N$, (since $K+I=N$ and $I\geq 1$)
\item $f_{I-1,J,K-2I+1}$ for which $(I-1)+(K-2I+1) = K - I < N$ (since $K+I=N$ and $I\geq 1$).
\end{enumerate}
This means that the RHS of \eqref{eqn:frel} has been determined already, and this determines
the LHS uniquely.
\end{proof}

